% chktex-file 18
\chapter{Weltraumgestützte Systeme}\label{ch:weltraumgestuetztesysteme}

Weltraumgestützte Systeme nutzen unterschiedliche Sensorplattformen, Orbit"=Konfigurationen und Instrumente, die je nach Missionsziel spezifische Aufgaben im Rahmen des Waldbrandmanagements übernehmen. Grundsätzlich operieren die meisten dieser Satelliten in erdnahen, oftmals sonnensynchronen Umlaufbahnen (Low Earth Orbit, LEO), um kontinuierlich Landoberflächen streifenweise abzutasten~\cite{TerraLexikon}.

Für die globale Erkennung von aktiven Waldbränden kommen zunehmend moderne Systeme wie die Visible Infrared Imaging Radiometer Suite (VIIRS) zum Einsatz. VIIRS nutzt Sensoren, die auf verschiedenen Satellitenplattformen installiert sind, darunter die Suomi National Polar-orbiting Partnership (Suomi NPP) und das Joint Polar Satellite System (JPSS). Es kombiniert die Aufgaben älterer Sensorgenerationen wie MODIS und AVHRR, die beispielsweise auf dem Terra-Satelliten oder auf der NOAA-POES- und MetOp-Serie installiert waren und bietet eine verbesserte räumliche Auflösung sowie höhere Detektionsgenauigkeit~\cite{MODISLexikon,TerraLexikon}.
VIIRS arbeitet als scannendes Radiometer, welches die Erdoberfläche im sichtbaren und infraroten Spektrum abtastet. Dadurch ist eine großflächige, weltweite und automatisierte Erfassung von Landoberflächentemperaturen und feuerbedingten Anomalien möglich~\cite{VIIRSLexikon}.

Zur Erfassung von Ereignissen mit hohen Temperaturen, wie sie unter anderem bei Waldbränden auftreten, werden Kleinsatelliten-Missionen genutzt. Diese sind spezifisch für den Zweck der Feuerfernerkundung konzipiert worden. Eine Pionierrolle wird hierbei von der deutsche BIRD-Satellit eingenommen, welche weltweit der erste seiner Art war. Auf diesen Erkenntnissen aufbauend wurde die FireBIRD-Mission mit den beiden Satelliten TET-1 (Technologie-Erprobungsträger) und BIROS (Bi-spectral Infrared Optical System) entwickelt. Sie ermöglichen nicht nur die reine Detektion, sondern können auch weitere wichtige Parameter, wie die Feuerintensität und die Charakterisierung der Feuerfront, mit höherer räumlicher Auflösung als vergleichbare Systeme liefern~\cite{DZLR2024}.
Ein zukunftsweisendes Konzept der Kleinsatelliten-Missionen bietet unter anderem das Unternehmen OroraTech, welches acht Satelliten in den Orbit gebracht hat. Die Anordnung dieser Satelliten bietet den entscheidenden Rahmen für zwei wesentliche Innovationen: Die höhere Auflösung des Systems ermöglicht die frühzeitige Erkennung von Bränden bereits ab einer Größe von 4\,$\times$\,4\,m. Dies gewährleistet eine zeitnahe Alarmierung der Einsatzkräfte nach Brandentstehung. Des Weiteren liefern die Satelliten Daten für die Nachmittags- sowie Nachtstunden. Diese Zeiträume wurden bisher von den meisten Systemen kaum abgedeckt.
 Somit lassen sich nun auch Phasen besonders intensiver Brandaktivität lückenlos überwachen. Zur Analyse der enormen Datenmengen und zur Minimierung der Latenzzeit werden die Daten mithilfe künstlicher Intelligenz ausgewertet. Durch das sogenannte Edge-Computing erfolgt die Datenauswertung meist direkt auf dem Satelliten vor der Datenübertragung, sodass die Ergebnisse in Echtzeit vorliegen können~\cite{Seton2025, DZLR2024}.


\chapter{Optisch-terrestrische Systeme}\label{ch:optischTerrestrischeSysteme}

Um den Herausforderungen der Waldbrandfrüherkennung und -überwachung gerecht zu werden, werden neben weltraumgestützten Systemen auch optisch-terrestrische Systeme eingesetzt. Diese bodengebundenen ("terrestrischen") Technologien basieren auf der Verwendung von visueller ("optischer") Kamerasensorik und sind stationär auf ehemaligen Feuerwachtürmen, Gebäuden oder anderen erhöhten Standorten installiert.
Moderne Systeme integrieren eine Sensoreinheit zur Bilddatenerfassung, eine Recheneinheit zur Datenverarbeitung sowie eine KI-basierte Analysesoftware. Neuronale Netze und Deep-Learning-Algorithmen ermöglichen die automatische Erkennung von Rauch- und Feuersignaturen, sodass zielgerichtet und zeitnah gehandelt werden kann~\cite{IQGmbH2023,Winter2021}.

Die Positionierung der Systeme erfolgt mit einer Mindesthöhe von fünf Metern oberhalb der Baumkronen, um ein ungehindertes Sichtfeld zu gewährleisten und somit die frühzeitige Erkennung von aufsteigenden Rauchsäulen zu ermöglichen.
Um eine lückenlose Erfassung zu gewährleisten, sind die Systeme in der Regel mit einem rotierenden Schwenk-Neige-Kopf ausgestattet. Dieser ermöglicht eine Überwachung des Gebietes in einem \qty{360}{\degree}-Radius. Das vom französischen Unternehmen Paratronic entwickelte Erkennungssystem ADELIE (Alert Detection Localisation of Forest Fire) benötigt beispielsweise für eine vollständige Umdrehung maximal zwei Minuten~\cite{Gemmingen2022}.

Hinsichtlich der Sensorik verwenden moderne Systeme hochauflösende, multispektrale Kameras, um unter unterschiedlichen Tages- und Witterungsverhältnissen eine zuverlässige Erkennung zu gewährleisten.
Das ursprünglich vom DLR (Deutsches Zentrum für Luft- und Raumfahrt) für Weltraummissionen entwickelte System "IQ FireWatch" integriert bis zu vier Sensoren in einem Gehäuse. Dazu gehören ein Monochrom-Sensor für die Aufnahme von Bildern am Tag, ein Nah-Infrarot-Sensor (NIR) für Aufnahmen bei Nacht sowie ein RGB-Sensor für zusätzliche Spektralinformationen. Optional kann, vor allem für Industriegelände, ein zusätzlicher thermischer Infrarotsensor integriert werden, der primär auf kleineren Entfernungen eine gesicherte Erkennung von Hitzequellen ermöglicht. Unter optimalen Bedingungen ist das System so in der Lage, Brände bis zu einer Entfernung von 60~km zu erkennen~\cite{IQGmbH2023, Winter2021}.

Die eigentliche Brandfrüherkennung und Automatisierung des Systems erfolgt über eine komplexe Software, die Rohdaten in Echtzeit auswertet. Das Erkennungssystem ADELIE registriert die Bilddaten und vergleicht sie auf Veränderungen, etwa die Verschiebung von Pixeln, die auf Rauch oder Feuer hindeuten könnten. Anschließend werden diese Differenzen unter Zunahme von verschiedenen Algorithmen analysiert, um alle Elemente, die keinen Rauch darstellen, bestmöglich zu eliminieren~\cite{Winter2021}.

Die Analyse moderner Systeme setzt dabei zunehmend auf einen hybriden Ansatz: Klassische, merkmalbasierte Verfahren werden durch Künstliche Intelligenz (KI) ergänzt. Die neuronalen Netzwerke der KI verarbeiten die mehrdimensionalen Bildmatrizen und erlernen komplexe Muster, um tatsächlichen Brandrauch eindeutig von anderen natürlichen Störfaktoren oder Bildveränderungen zu unterscheiden~\cite{IQGmbH2023,Winter2021}.

Wird auf diesem Weg potenzieller Rauch erkannt, schätzt das System die Entfernung zum potenziellen Brandherd ab. Zudem wird geprüft, ob sich der Rauch in einem Bereich befindet, in dem bekannte Rauchquellen, wie Industrieanlagen und Schornsteine, existieren. Dadurch können Fehlalarme minimiert werden. Das System alarmiert die Einsatzkräfte danach jedoch nicht automatisch, sondern übermittelt die Informationen zunächst an die zuständige Leitstelle, damit diese durch einen menschlichen Operator validiert werden können. Die Mitarbeiter nehmen daraufhin selbst mithilfe einer hochauflösenden Kamera mit starkem optischem Zoom eine Prüfung vor, um festzustellen, ob es sich faktisch um einen Brand handelt, bevor sie die Einsatzkräfte alarmieren.
Das Gesamtsystem agiert folglich als intelligentes Entscheidungsunterstützungssystem~\cite{IQGmbH2023, Gemmingen2022}.


\chapter{IoT-basierte Sensornetzwerke}\label{ch:IoTbasierteSensornetzwerke}

Während satelliten- und optisch-terrestrische Systeme oft erst nach der Entstehung einer offenen Flamme erkennen, setzten IoT-basierte Sensornetzwerke direkt auf der Ebene des Waldbodens und damit deutlich früher an. Brände können bereits in der Schwelbrandphase erkannt werden, bevor sich sichtbare Flammen oder Rauchwolken entwickeln. Dafür greift die Technik auf \enquote{Digitale Nasen} (Electronic Noses oder e-Noses) zurück, die versuchen, das biologische Geruchssystem von Menschen und Tieren nachzuahmen, um spezifische Gerüche und luftgetragene chemische Verbindungen zu erkennen und zu identifizieren.
Zu den grundlegenden Komponenten einer digitalen Nase gehören in der Regel Arrays aus chemischen Sensoren, Signalverarbeitungseinheiten und Mustererkennungsalgorithmen. Nimmt eine digitale Nase Brandgase wie Kohlenstoffmonoxid, Wasserstoff oder flüchtige organische Verbindungen (VOCs) wahr, reagieren Sensoren auf diese Stoffe und wandeln die chemischen Informationen in elektrische Signale um. Diese Signale bilden ein komplexes, mehrdimensionales Geruchsprofil, welches von Algorithmen des maschinellen Lernens und künstlichen Intelligenz mit bekannten Profilen von Feuer und Rauch abgeglichen wird~\cite{Chan2024, FayosJordan2024, Kumar2023}.

Ein wegweisendes Beispiel ist das laufende Forschungsprojekt KIWB (KI-basiertes Warnsystem zur Waldbranderkennung). Es setzt bereits in der Entstehung an und hat zum Ziel, Schwelbrände bereits ab einer Ausdehnung von nur einem Quadratmeter sicher zu detektieren. Das System basiert auf dem Einsatz von energieautarken Gassensoren, die über weite Distanzen funkfähig und für den Langzeiteinsatz im Wald ausgelegt sind. Mittels spezialisierter KI-Lösungen sollen die Sensoren spezifische Gassignaturen von Bränden erkennen und von alltäglichen Emissionen unterscheiden. So soll eine zuverlässige Früherkennung ermöglicht werden, die gleichzeitig die Senkung auf eine Fehlalarmquote von unter ein Prozent anstrebt~\cite{BBF2025}.

Im kommerziellen Sektor setzt das Unternehmen Dryad Networks diese Technologie mit seinem Silvanet-System großflächig ein. Das System besteht aus solarbetriebenen, KI-gestützten Gassensoren, die direkt an den Bäumen des Waldes installiert werden und Brände innerhalb weniger Minuten nach der Entzündung erkennen können. Zur Übertragung der Alarmsignale aus unwegsamen Waldgebieten an die Leistelle, wird eine spezielle Funkinfrastruktur genutzt, da klassische Mobilfunknetzte dort oft nicht oder nur schlecht zur Verfügung stehen~\cite{DNG}.

Hierfür kommt das sogenannte Low-Power Wide-Area Network (LPWAN) zum Einsatz. Diese Technologie ist darauf ausgerichtet kleine Datenmengen extrem energieeffizient über große Distanzen zu senden. Der am weitesten verbreitete und lizenzfreie Standard für LPWAN ist LoRaWAN (Long Range Wide Area Network). LoRaWAN ist ein weltweit offenes Kommunikationsprotokoll speziell für energieeffiziente IoT-Geräte und findet Einsatz in den verschiedensten Anwendungsbereichen, von der Landwirtschaft bis hin zur städtischen Infrastruktur. Es ermöglicht die Übertragung von Daten über Entfernungen von mehreren Kilometern, selbst in dicht bewaldeten oder bergigen Gebieten.
Innerhalb dieses Netzwerks übernehmen sogenannte Gateways eine zentrale Rolle. Anstatt direkt mit einem weit entfernten Server zu kommunizieren, senden die Sensorknoten ihre Pakete zunächst an ein nahegelegenes Gateway, welches die empfangenen Daten dann an einen mit fortschrittlicher Technologie ausgestatteten Server weiterleitet. Klassische LoRaWAN-Systeme sind dabei meist sternförmig aufbaut, es gibt jedoch auch Ausnahmen. Das System Silvanet von Dryad Networks nutzt beispielsweise eine Mesh-Topologie, bei der die Daten von den Gateways untereinander weitergereicht werden~\cite{Nodes.sh2024}. 

durch die Kombination aus direkter Gasmessung am Boden, lokaler Datenverarbeitung durch KI und energieeffizienter Funkübertragung bilden IoT-Sensornetzwerke ein wichtiges Element in der Waldbrandfrüherkennung. Sie schließen die Überwachungslücke im kritischen Anfangsstadium eines Brandes und ermöglichen damit eine noch schnellere Reaktion~\cite{Chan2024,DNG}.
